\begin{problem}
  Prove that the following equations are identities:

  \begin{enumerate}
    \item $\sec (t) - \cos (t) = \tan (t)\sin (t)$ [$3.75$ points]
    \item $\frac{1}{1 - \sin (t)} + \frac{1}{1 + \sin (t)} = 2\sec^{2} (t)$ [$3.75$ points]
  \end{enumerate}
\end{problem}

\begin{solution}
  \begin{enumerate}
    \item $\sec (t) - \cos (t) = \tan (t)\sin (t)$ [$3.75$ points]

      \begin{align*}
        \sec (t) - \cos (t) &= \tan (t)\sin (t) \\
        \frac{1}{\cos (t)} - \cos (t) &= \frac{\sin (t)}{\cos (t)} \times \sin (t) \\
        \frac{1}{\cos (t)} - \cos (t) &= \frac{\sin^{2} (t)}{\cos (t)} \\
        \cos (t) \left(\frac{1}{\cos (t)} - \cos (t)\right) &= \left(\frac{\sin^{2} (t)}{\cos (t)} \right) \\
        1 - \cos^{2} (t) = \sin^{2} (t) \\
      .\end{align*}

    \item $\frac{1}{1 - \sin (t)} + \frac{1}{1 + \sin (t)} = 2\sec^{2} (t)$ [$3.75$ points]

    \begin{align*}
      \frac{1}{1- \sin (t)} + \frac{1}{1 + \sin (t)} &= 2\sec^{2} (t) \\
      \frac{(1 + \sin (t))(1 - \sin (t))}{(1 - \sin (t))(1 + \sin (t))} &= 2\sec^{2} (t) \\
      \frac{2}{1 - \sin^{2} (t)} &= 2\sec^{2} (t) \\
      2 \left(\frac{1}{1 - \sin^{2} (t)}\right) &= 2\sec^{2} (t) \\
      \frac{1}{1 - \sin^{2} (t)} &= \sec^{2} (t) \\
      \frac{1}{\sec^{2} (t)} &= 1 - \sin^{2} (t) \\
      \cos^{2} (t) &= 1 - \sin^{2} (t) \\
    .\end{align*}
  \end{enumerate}
\end{solution}

\newpage

\begin{problem}
  Suppose that $\frac{\pi}{2} < A < \pi$ and $\sin (A) = \frac{\sqrt{13}}{7}$
  and $\cos (A) = -\frac{6}{7}$. Use appropriate identities to find
  \textbf{exact} values for the following. (Do \textit{not} use an inverse
  trig function to solve or $A$.)

  \begin{enumerate}
    \item $\sin (2A)$ [$2.5$ points]
    \item $\cos (2A)$ [$2.5$ points]
    \item $\cos (\frac{2}{A})$ [$2.5$ points]
  \end{enumerate}
\end{problem}

\begin{solution}
  \begin{enumerate}
    \item $\sin (2A)$ [$2.5$ points]

      To solve this, we can use the Double-Angle Identities for sin:

      \begin{align*}
        \sin (2\theta) &= 2\sin(\theta)\cos(\theta) \\
                       &= 2 \left(\frac{\sqrt{13}}{7}\right) \left(-\frac{6}{7}\right) \\
                       &= \frac{-12\sqrt{13}}{49}
      .\end{align*}

    \item $\cos (2A)$ [$2.5$ points]

      To solve this, we can use any of the $3$ Double-Angle Identities for cos:

      \begin{align*}
        \cos (2\theta) &= 2\cos^{2} (\theta) - 1 \\
                       &= 2\left(-\frac{6}{7}\right)^{2} - 1 \\
                       &= 2\left(\frac{36}{49}\right) - 1 \\
                       &= \frac{72}{49}
      .\end{align*}

    \item $\cos (\frac{2}{A})$ [$2.5$ points]

      To solve this, we can use the Half-Angle Identity for cos.

      \begin{align*}
        \cos (\frac{\theta}{2}) &= \pm \sqrt{\frac{1 + \cos (\theta)}{2}} \\
                                &= \pm\sqrt{\frac{1 + -\frac{6}{7}}{2}} \\
                                &= \pm\sqrt{\frac{\frac{1}{7}}{2}} \\
                                &= \pm\frac{\sqrt{14}}{14}
      .\end{align*}
  \end{enumerate}
\end{solution}
